% -----------------------------------------------------------------------------------------------------------------------------------------
% 第9章 結論：都市計画への展開

% -----------------------------------------------------------------------------------------------------------------------------------------

\chapter{結論：都市計画への展開}
\label{chap:conclusion}

\begin{center}
\fbox{\textbf{【編集中】}}
\end{center}

\vspace{1em}

% -----------------------------------------------------------------------------------------------------------------------------------------
\section{研究の総括}
\label{sec:ch7_summary}

% -----------------------------------------------------------------------------------------------------------------------------------------
\subsection{研究目的の達成}

本研究は3つの目的を掲げた。公共交通政策における実装ギャップの要因として人間の認知バイアスの影響を計算論的に解明すること、生成AIと人間の協調的関係性を具体化するシステムとして秘密保護型政策評価システムの可能性を探ること、そして認知バイアスと生成AIを考慮した制度設計への示唆を導出することである。

\paragraph{目的1について}
第5章において、協調ロボット制御モデルを用いた計算論的分析により、三つの認知バイアスが政策協調に与える影響を定量的に解明した。特に、現状維持バイアスの閾値効果、確証バイアスの逆説的効果、狭い視野の一貫した負の影響を発見した。

\paragraph{目的2について}
第6章において、秘密保護と民主的正当性を両立する政策評価システムを提案し、LLM as a Judge、Constitutional AI、市民討議を組み合わせたアーキテクチャを設計した。

\paragraph{目的3について}
第9章において、認知バイアスへの対処戦略、生成AIを組み込んだ制度設計、「連携・共創」の再設計に向けた枠組みを提示した。

% -----------------------------------------------------------------------------------------------------------------------------------------
\subsection{核となる主張}

本研究の核となる主張は、以下の通りである。

\begin{quote}
\textbf{生成AIは人間の「執政の創造性」を補完する「杖」として、認知バイアスへの対話的介入を通じて政策形成を支援できる。秘密保護型評価の仕組みを活用することで、秘匿性と信頼性を両立した政策評価システムが可能になる。}
\end{quote}

% -----------------------------------------------------------------------------------------------------------------------------------------
\section{理論的貢献}
\label{sec:ch7_theoretical_contributions}

% -----------------------------------------------------------------------------------------------------------------------------------------
\subsection{認知バイアスの計算論的分析手法の政策科学への導入}

本研究は、協調ロボット制御理論を政策ネットワーク分析に応用し、認知バイアスの影響を計算論的に解明する手法を導入した。このアプローチは、政策科学における計算論的転回（computational turn）の一環として位置づけられる。

% -----------------------------------------------------------------------------------------------------------------------------------------
\subsection{秘密保護型政策評価の設計}

本研究は、秘匿証明の考え方を政策評価に応用し、「秘密を守りながら専門性を証明する」という新たな政策参加のあり方を設計した点で先駆的な試みである。% [校正メモ]: 先行研究との差異を具体的に確認・記述する（Task \#10）

% -----------------------------------------------------------------------------------------------------------------------------------------
\subsection{生成AIの「杖」としての理論的位置づけ}

本研究は、生成AIと人間の関係性を「杖」として補完的に位置づける理論的枠組みを提示した。これは、AIによる代替ではなく、AIと人間の協調を前提とする視点である。

% -----------------------------------------------------------------------------------------------------------------------------------------
\section{実践的貢献}
\label{sec:ch7_practical_contributions}

% -----------------------------------------------------------------------------------------------------------------------------------------
\subsection{制度設計への提言}

本研究は、認知バイアスと生成AIを考慮した制度設計への具体的な提言を行った。現状維持バイアスへの対処として、段階的変化導入の設計を提案した。確証バイアスへの対処として、「悪魔の代理人」の制度化を提案した。狭い視野への対処として、横断的評価指標の導入を提案した。生成AIの活用として、補完性、透明性、アカウンタビリティの原則を提示した。

% -----------------------------------------------------------------------------------------------------------------------------------------
\subsection{秘密保護型政策評価システムの設計指針}

本研究は、秘密保護型政策評価システムの具体的な設計指針を提示した。三層アーキテクチャの採用、Constitutional AIと市民討議による評価基準設計、控訴プロセスによる人間介入の確保がその主要な内容である。

% -----------------------------------------------------------------------------------------------------------------------------------------
\section{歴史的視座：日本における合理的政策転換の試み}
\label{sec:ch7_historical}

% -----------------------------------------------------------------------------------------------------------------------------------------
\subsection{EBPMと歴史的先行事例}

近年、EBPM（証拠に基づく政策形成）が注目を集めているが、日本において合理的な政策への転換を提唱する試みは決して新しいものではない。1934年、陸軍の永田鉄山は『国防の本義と其強化の提唱』を著し、従来の主観的・慣習的な政策決定から、科学的根拠に基づく合理的な政策への転換を強く主張した\footnote{永田鉄山「国防の本義と其強化の提唱」1934年。本書は、総力戦体制の構築を念頭に、国家資源の科学的な動員と管理の必要性を説いた。}。

永田は、政策決定において「科学的調査」と「合理的判断」を重視し、感情論や慣習に基づく決定を批判した。この主張は、今日のEBPMが目指す「証拠に基づく意思決定」と理念的に共通するものであった。しかし、当時の政治的文脈において、この合理的アプローチは十分に実現されなかった。軍部内の派閥対立や、政治的圧力により、科学的根拠に基づく政策形成は阻害されたのである。

この歴史的事実は、本研究にとって重要な示唆を与える。合理的政策形成の理念そのものは技術の進歩とは独立して存在し得るが、その実現には技術的条件だけでなく制度的・政治的条件が不可欠であり、認知バイアスや政治的圧力がいかに強力な阻害要因になり得るかをこの事例は示している。

% -----------------------------------------------------------------------------------------------------------------------------------------
\subsection{生成AI時代における新たな可能性}

永田の時代と比較して、今日の状況は本質的に異なる。データの収集・分析能力が飛躍的に向上し、生成AIという新たな技術的基盤が登場し、市民社会の成熟と民主的制度の定着が進んだ。

しかし、認知バイアスという人間の根本的な特性は変わらない。本研究が明らかにしたように、現状維持バイアスや狭い視野は、協調的ガバナンスを阻害し続ける。永田の時代には技術的に解決できなかったこれらの課題に対し、生成AIを「杖」として活用することで、新たな解決の道が開かれる可能性がある。

もちろん、歴史の教訓を忘れてはならない。技術的な解決策が存在しても、制度的・政治的な障壁を克服できなければ、理念は実現しない。本研究が提案する秘密保護型政策評価システムも、技術的可能性にとどまらず、制度的設計と社会的受容を含む包括的なアプローチとして位置づけられるべきである。


\section{今後の課題：都市計画を舞台にした実証}
\label{sec:ch7_future_work}

% -----------------------------------------------------------------------------------------------------------------------------------------
\subsection{より複雑な政策領域への適用}

公共交通政策は、本研究の「実験場」として適切であったが、より複雑な政策領域への適用が期待される。特に、都市計画は多様なステークホルダー（土地所有者、開発業者、住民、行政など）、多様な秘密情報（土地利用計画、開発権、資産価値など）、長期的影響（数十年単位での都市構造の変化）という点で興味深い研究対象となる。

% -----------------------------------------------------------------------------------------------------------------------------------------
\subsection{秘密保護型政策評価システムの社会実装}

秘密保護型政策評価システムの社会実装に向けては、技術的・制度的・社会的な課題が残されている。技術的にはTEE、LLM、Constitutional AIの統合が必要であり、制度的には法的位置づけや運用ルールの策定が求められ、社会的には市民の理解と信頼の獲得が重要である。

% -----------------------------------------------------------------------------------------------------------------------------------------
\subsection{生成AIと人間の協調的関係性の継続的検証}

生成AI技術は急速に進化しており、人間-AI協調のあり方も変化し続ける。本研究の枠組みは、継続的な検証と改善を必要とする。

% -----------------------------------------------------------------------------------------------------------------------------------------
\section{結び}
\label{sec:ch7_conclusion}

本研究は、公共交通政策を舞台に、生成AIと人間の協調的関係性を探求した。その過程で、認知バイアスが政策協調に与える異質な影響を計算論的に識別し、秘密保護と民主的正当性を両立する政策評価システムを提案し、制度設計への示唆を導出した。

公共交通政策は、本研究の「第一の舞台」であった。今後は、都市計画という「第二の舞台」での実証を通じて、生成AIと人間のより良い協調のあり方を探求していきたい。

% -----------------------------------------------------------------------------------------------------------------------------------------


% -----------------------------------------------------------------------------------------------------------------------------------------
\printbibliography[heading=subbibliography,title={章末参考文献}]
